%%%%% Section 7. Discussion %%%%%
% 2026-08-03 全面差し替え。旧稿(Ver.2 コピー)は破棄。
% §7.1 のみ執筆済み。残る subsections(図書館への含意/予期される反論/
% 表現とインターフェース 等)は未執筆(instructions.md §4 参照)。
% 2026-08-08: 応用システム/rating-as-content/record への追記(客観性・普遍性)の
% 3 subsections を末尾に追加(見出し仮)。内容メモ(コメント)+暫定散文(Claude 叩き台)。
\section{Discussion}\label{sec:discussion}

%%% Subsection 7.1 Where Connotation Lives %%%%%
% 見出しは仮(2026-08-03)
\subsection{Where Connotation Lives}\label{subsec:where_connotation}

% --- P1: sec2 の record の非対称 → 単位の非対称へ ---
Section~\ref{subsec:naming} located the difference between the two
kinds of attribute at the record: denotation arrives with it,
connotation is written in no field. The bundle locates the
difference more deeply---at the unit. A name is attached to one
work at a time; what the OT-FB expresses is not a property of one
work at all.

% --- P2: 三段の認識論(pairwise で不可視 → ドメイン内で構造 → ドメイン間で independence)---
Consider a collection of two works. Between them there is a
distance, and nothing more---nothing to say how much of it is
genre and how much runs across genre. Connotation is not hidden
in such a world; it has nowhere to exist. Within a genre, among
many works, it takes shape: each work shares some connotations
with some works and not with others, and this web of shared and
unshared gives every work a relative position among its own kind.
And it is only across genres that the last question can be put:
which part of this shape is independent of the genre itself?
Where the ensembles correspond as wholes, galaxy laid onto
galaxy, what the correspondence preserves is shared across
domains---and being shared across domains is what
\emph{independent of denotation} means. Such a whole-to-whole
correspondence is what transport computes, and what no pairwise
comparison can.

% --- P3: だから single word にならない(語彙ではなく単位の不一致)---
This is why connotation was never going to be carried whole by a
name. A name attaches to the individual work; connotation, as the
bundle draws it, lives in the relations among domains. The
consistent failure of connotative naming that
Section~\ref{subsec:naming} recorded is, on this reading, not a
poverty of vocabulary but a mismatch of unit---and a coordinate
does not enrich the vocabulary; it changes the unit.

% --- P4: principal connotation =道しるべ(切り捨ては設計判断)---
Nor does the bundle keep every connotation. The principal
connotation analysis of Section~\ref{subsec:bundle_reach} retains
what is independent of denotation, present in the content data, and
principal among what remains; whatever depends on the genre is
left with the genre. This is a design decision, not a shortfall.
An axis whose meaning shifted from floor to floor could not guide
a leap; what can serve the explorer as a signpost is exactly what
keeps its meaning on every floor. The rabbit hole is dependable
for the same reason it is partial.

%%% Subsection 7.2(仮)Reading Connotation from Language %%%%%
% 2026-08-11 新設。内容メモ(コメント)+暫定散文(Claude 叩き台)。見出し・位置仮。
% 位置の根拠: §7.1 が「connotation はどこに宿るか」(関係・単位)を言い、本節が
% 「その関係をどのデータから、誰が読むのか」を引き受け、§7.x(From Representation
% to Systems)の query by description へ渡す。したがって §7.1 の直後・§7.x の直前。
%
% --- 内容メモ(2026-08-11 せんせーとの対話で確定)---
% (0) 本節が回収する伏線: §2.2 の `An information system cannot rely on such
%     background implicitly`(shared background = language, culture, embodied
%     experience。Clark1996 / Langacker2008 / Hjorland1995400)。
%     回答は「背景に頼らない」ではなく「背景の痕跡から読む読み手をシステムが所有する」。
% (1) **成立条件としての generic feature vector**(せんせー 2026-08-11)。
%     かつて feature vector は用途ごとに設計するものだった——何を特徴とするかを
%     設計者が事前に決める。それは名前を事前に選ぶことであり、§2.2 の Osgood の
%     `the scales themselves must still be chosen in advance` と同じ構造。
%     すなわち denotative manner で connotation を作ることと五十歩百歩だった。
%     generic embedding は用途に先立って存在し、設計者は次元の意味を選ばない。
%     だから初めて「データが軸を決める」余地が生じた。**これが本手法の背景**。
% (2) LLM は言語化装置ではなく**読者**。タグ列は言語(denotative なラベル列)だが、
%     connotation は列そのものではなく「列を読んだときに立ち上がるもの」に宿る。
%     ベクトルはその読みの残したもの。
% (3) 主張の限定——存在論ではなく操作論(せんせー 2026-08-11 確定)。
%     ×「LLM は embodiment を備える」/ ×「meaning を再現する」
%     ○「間接的かつ部分的ではあれ、connotative な情報が内部のベクトル表現から
%        引き出せる」。検証可能な主張の形に置き直すことで、決着しない存在論を通らない。
%     「部分的」は新たな譲歩ではなく §4.3 principal への三つの限定と同じもの
%     (→ §7.1 末尾 `dependable for the same reason it is partial` と韻を踏む)。
% (4) seeker 側と work 側は同じ読み手。LLM は現にプロンプトから connotation を
%     読む(読者自身の日常経験)。同じ読み手が works の側も読むから、query と work が
%     同じ空間に落ちる——§7.x の query by description の**機構がここで与えられる**
%     (現行の暫定散文 `lands somewhere in the connotation space` は機構を語っていない)。
%     ただし「現に読み解いている」は振る舞いの証拠であって表現の幾何の証拠ではない。
%     ゆえに順序は「幾何を測った研究(P3)→ 日常経験(P4)」とし、逆にしない。
% (5) 反対側への態度: form から meaning へ到達するとは主張しない。得るのは works の
%     **配置**であり、座標は何も名指さない(sec3)。BenderKoller2020 / Bisk2020 の
%     否定命題を認めたまま本論文の主張は無傷で通る。敬意を持って引く。
%
%%% TODO 引用(2026-08-11 Web 検索で実在確認済み。bibkey は仮・**ページ未確認**——
%%% ref.bib 追加時に一次ソースで巻号ページを検証すること。命名は既存則(著者+年+開始頁)に揃える):
%%%  (i)  Louwerse & Zwaan (2009) "Language Encodes Geographical Information,"
%%%       Cognitive Science 33(1). 都市名の共起統計のみから緯度経度が復元できる。
%%%       「言語には誰も書かなかったものが刻まれている」の最も直截な実証。→ 主柱1
%%%  (ii) Louwerse (2011) "Symbol Interdependency in Symbolic and Embodied Cognition,"
%%%       Topics in Cognitive Science 3(2). symbol interdependency hypothesis——
%%%       言語のシンボル体系は身体的・知覚的経験の痕跡を符号化している。
%%%       せんせーの「間接的に身体性を備えている」に認知科学側で付いている名前。→ 主柱2
%%%       (関連: Louwerse (2018) "Knowing the Meaning of a Word by the Linguistic
%%%        and Perceptual Company It Keeps," Topics in Cognitive Science 10(3).)
%%% (iii) Grand, Blank, Pereira & Fedorenko (2022) "Semantic projection recovers rich
%%%       human knowledge of multiple object features from word embeddings,"
%%%       Nature Human Behaviour 6(7). word embedding から**軸への射影**で人間の
%%%       特徴知識(size/danger/intelligence 等)を復元。§4.3 principal connotation
%%%       analysis と操作が同型。掲載誌の格も JDoc 読者に効く。→ 主柱3(最重要)
%%%  (iv) Marjieh et al. (2024) "Large language models predict human sensory judgments
%%%       across six modalities," Scientific Reports 14(arXiv:2302.01308). 補強。
%%%  (v)  Abdou et al. (2021) "Can Language Models Encode Perceptual Structure Without
%%%       Grounding? A Case Study in Color," CoNLL. 色の知覚構造。補強。
%%%  (vi) 反対側(敬意を持って引く): Bender & Koller (2020) "Climbing towards NLU,"
%%%       ACL / Bisk et al. (2020) "Experience Grounds Language," EMNLP.
%%% 既に ref.bib にあり本節で再出させうるもの: Clark1996, Langacker2008,
%%% Hjorland1995400(§2.2 で既出), Deerwester(LSA, 未使用)。
\subsection{Reading Connotation from Language}\label{subsec:reading_connotation}

% --- 以下、メモを実現する暫定散文(2026-08-11、Claude 叩き台)---
% --- P1: 成立条件(generic feature vector)= 本手法の背景 ---
The method has a precondition, and it is recent. Feature vectors were
once built for a purpose: whoever built them decided beforehand what
would count as a feature, and that decision was a choice of names.
Osgood's scales had to be settled before the rating began
(Section~\ref{subsec:naming}); a hand-built feature set makes the same
choice, in numbers rather than in words, and arrives at connotation by
denotative means. Generic embeddings come before any purpose. Nobody
has decided, for this collection, what the dimensions are to mean---
nothing in them has been named. Only under that condition is there room
left for the data to decide where the axes lie.

% --- P2: LLM は読者である(§2.2 の伏線回収)---
What such an embedding reads is language, and yet connotation is not in
the language it reads. A list of tags---genre, mood, era,
instrument---is a string of denotative labels; read them together, and
something rises that no one of them holds. This is the work the language
model is put to. It is not asked to name, but to read, and the vector is
what its reading leaves behind. Section~\ref{subsec:naming} observed
that between people connotation travels unnamed, on a shared background
of language, culture, and embodied
experience~\cite{Clark1996,Langacker2008,Hjorland1995400}, and that a
system cannot lean on that background implicitly. It need not. It can
hold a reader who has come by that background at second hand.

% --- P3: 主張の限定 + 認知科学の足場(測られていること)---
How much of the background such a reader has is a fair question, and the
claim we make is narrow. Not that a language model has embodied
experience; not that it recovers meaning. Only that connotative
information is present in its internal vectors and can be drawn out of
them---indirectly, and in part. That much has been measured rather than
assumed. Language carries the trace of perceptual experience: from the
co-occurrence statistics of city names alone, latitude and longitude can
be recovered~\cite{LouwerseZwaan2009}---a regularity Louwerse names
symbol interdependency~\cite{Louwerse2011}. Human judgments of object
features are recovered from word embeddings by projection onto an
axis~\cite{Grand2022}, which is the operation of
Section~\ref{subsec:bundle_reach} performed on other data. Language
models predict human sensory judgments across
modalities~\cite{Marjieh2024}, the structure of color among
them~\cite{Abdou2021}. Whether form can suffice for meaning remains
contested~\cite{BenderKoller2020,Bisk2020}, and this construction does
not need it to: what it claims to obtain is an arrangement of works, not
their meaning, and the coordinate marks a place while naming nothing
(Section~\ref{sec:framework}).

% --- P4: 一人の読み手が両側を読む → §7.x の query by description へ渡す ---
One reader serves both sides. A language model reads connotation from a
prompt as readily as from a list of tags---which is the ordinary
experience of anyone who has used one. The groping description of
Section~\ref{subsec:representation_to_systems}, then, and the works it
gropes for are read by the same reader, and fall into the same space.
The seeker's few words and the collection's many are commensurable
because of who reads them.


%%% Subsection 7.x(仮)From Representation to Systems %%%%%
% 2026-08-08 下書き。内容メモ(コメント)+暫定散文(Claude 叩き台)。
% 見出し・位置・順序はすべて仮。既定の残り subsections
% (図書館への含意/予期される反論/表現とインターフェース)との統合・分割は
% 執筆時に判断する。とくに「表現とインターフェース」(v9 §9)とは対句の関係:
% v9 §9 =「interface ではなく representation を」(貢献の位置づけ)、
% 本節=「representation が得られたいま、その上に interfaces は増殖する」(展望)。
\subsection{From Representation to Systems}\label{subsec:representation_to_systems}

% --- 内容メモ: disentangled representation からどんな情報システムが作れるか ---
% (1) ウサギ穴付き図書館: representation をそのまま探索空間にする。
%     従来型の denotative なシステムにウサギ穴を「追加」する形式も可能だし、
%     denotative space ごと探索可能な空間にすることもできる。
% (2) 関連作品による connotation 検索: 該当しそうな works を利用者が例示して
%     (例: Miles Davis の "Kind of Blue")、「connotation が共通する他の作品」を
%     検索・探索する。
% (3) Fuzzy description 検索: 「静かで、色彩的で、しかもそれらが溶け合ったような……」
%     という自然言語的表現や、関連しそうな単語列から connotation 検索・探索する。
% (4) 典型的タグによるスライダー検索(2026-08-08 追加): connotative axis と
%     相互情報量の高いタグへのマッチ度をスライダー(0/1 の二値も可)で指定して
%     検索・探索する。これを備えると Book House(アイコン入力)も取り込める。

% --- 以下、メモを実現する暫定散文(2026-08-08、Claude 叩き台)---
Section~\ref{sec:introduction} argued that representation must come
first, because no interface can surface what the space beneath it
does not encode. The converse deserves a word: once the space
exists, interfaces multiply. The most direct is the library of
Section~\ref{sec:framework} itself---rabbit holes added to an
otherwise conventional denotative catalog, or a space in which the
floors themselves are open to walking. A second is query by works:
a seeker offers what can be named---\textit{Kind of Blue}---and asks
for works that share its connotation, wherever their genres lie. A
third is query by description: a groping, natural-language
gesture---``quiet, and full of color, and the colors run into one
another''---lands somewhere in the connotation space, and the
neighborhood does the rest. A fourth returns an earlier idea to
service. Tags that share high mutual information with a connotative
axis can stand for that axis, and a seeker can set, slider by
slider, how much of each they want---continuously, or simply on and
off. This is recognizably the Book House~\cite{Pejtersen1983230},
its icons refounded: what once lived in the session can now read
from a space that outlives it. Different entrances, one space; what
varies is only how the seeker points.

%%% Subsection 7.y(仮)Rating Data as Content Data %%%%%
% 2026-08-08 下書き。内容メモ(コメント)+暫定散文(Claude 叩き台)。
\subsection{Rating Data as Content Data}\label{subsec:rating_as_content}

% --- 内容メモ: collaborative filtering の user-item rating data を content data と
%     みなして disentangled representation を作る ---
% * ウサギ穴(disentangle された axis)との相互情報量が高い curator には、
%   その名前をウサギ穴に付けることができる——「○○さんのウサギ穴」。
%   curator 全員ではなく、axis との関連度(相互情報量)の高さで基準となる
%   curator を見つける、という選び方。
% * ユーザー側から見れば、「自分の中にある何か」と resonate する curator を
%   探索的に見つけられる、という読み替えにもなる。
% * "curator" という語が適切かどうかは要検討(暫定散文では curator を仮使用)。
% * 既決との接続(2026-08-04 整理): rating data は content data と同等に扱える。
%   recsys との本質差は data source ではなく「ユーザーをモデル化するか」——
%   ここでは同じ rating data から works をモデル化する。§2.3 の書き方と揃える。

% --- 以下、メモを実現する暫定散文(2026-08-08、Claude 叩き台)---
Nothing in the construction requires that the content data be sound
or text. A work is described by those who attend to it as well as
by what it contains: the ratings a work gathers across many
listeners are themselves data about the work, and from such data
the same bundle can be built. This is the data of collaborative
filtering; but the difference between the two, as
Section~\ref{subsec:systems} drew it, lies not in the data source
but in what is modeled. Collaborative filtering models the user,
and searches on their behalf. The bundle built from rating data
still models the works, and leaves the walking to the seeker.

One consequence is worth drawing out. Among those whose choices
enter such data are curators---makers of playlists, keepers of
small collections. Where one curator's selections share high mutual
information with a connotative axis, the axis can carry the
curator's name: so-and-so's rabbit hole. The name does not define
the hole; the geometry has already drawn it, and naming, as ever,
presupposes representation. But for the seeker who descends it, the
name adds something no coordinate carries---a person, met only
through the works, whose way of hearing resonates with their own.

%%% Subsection 7.z(仮)What the Record Should Carry %%%%%
% 2026-08-08 下書き、同日改訂(せんせーメモ第2版に基づく)。
% 内容メモ(コメント)+暫定散文(Claude 叩き台)。見出し仮。
\subsection{What the Record Should Carry}\label{subsec:what_record_carries}

% --- 内容メモ(2026-08-08 改訂版)---
% * 「すべての connotation を集めた space」があるなら、おそらく無限次元。
%   connotation とは、人びとが生きて経験しているその瞬間ごとに、常に新しく
%   立ち上がるものだから(★要引用: connotation とはそういうものだと述べた論文。
%   おそらく LIS ではなく認知言語学の領域)。
% * 本論文が提供するのは、その一部を切り出した "principal connotation space"。
%   個々の work はこの部分空間の座標として位置づけられる。
%   → では、この座標系はどれだけ客観的か?
% * それは「何をもって principal な connotation とみなしたいか」という、
%   その図書館・情報アーカイブの運用目的が定める。あるいは音楽配信サイトや
%   電子ブックショップの経営方針が定める。各組織は入手可能な content data を用い、
%   目的に合ったメトリックを用いる——いわばメタレベルでの inductive bias。
% * 目的を共有する組織同士なら、互いに互換性のある connotation coordinate を
%   採用できる。それは record に次の項目を追加することに相当する:
%   「局所ユークリッド距離、あるいは標準内積が有効な content data そのもの」。
% * (前版メモからの継承、散文に残すかは要判断: 「客観的な基準が生じた瞬間に、
%   connotation は denotation になる」——暫定散文では客観性の問いへの入り口として
%   1文だけ残した。)

% --- 以下、メモを実現する暫定散文(2026-08-08 改訂、Claude 叩き台)---
If there were a space of all connotations, it would have no finite
dimension. Connotation is not a stock of properties awaiting their
index; it rises anew in every moment of lived experience, as often
as someone hears, reads, remembers.
%%% TODO 引用。「意味・概念は固定された在庫ではなく、使用・経験の瞬間ごとに
%%% 新しく構築される」を述べる文献の候補(2026-08-08 Web 検索で実在確認済み。
%%% どれを実際に引くかの判定と一次ソース検証は sec7 執筆時に行う):
%%% (a) Casasanto & Lupyan (2015) "All Concepts Are Ad Hoc Concepts," in Margolis &
%%%     Laurence (eds.), The Conceptual Mind, MIT Press. 主張が最も直截——概念は
%%%     使用のたびにその場で新しく構築される。
%%% (b) Croft & Cruse (2004) Cognitive Linguistics, Cambridge UP, Part II の
%%%     dynamic construal approach——語の意味は固定的でなく、使用ごとに on-line で
%%%     construal される。認知言語学の教科書的典拠として最も安全。
%%% (c) Barsalou (1987) "The instability of graded structure," in Neisser (ed.),
%%%     Concepts and Conceptual Development, Cambridge UP. 実験的基盤——同一人物でも
%%%     機会ごとに概念構造が変動する。前段に Barsalou (1983) "Ad hoc categories,"
%%%     Memory & Cognition 11.
%%% (d) Langacker (2008) Cognitive Grammar: A Basic Introduction, Oxford UP.
%%%     meaning is conceptualization——意味は処理時間の中で展開する動的な概念化。
%%% (e) 語用論側(記憶ベース・Web 未確認): Carston (2002) Thoughts and Utterances,
%%%     relevance theory の ad hoc concepts——語義は発話ごとに調整・構築される。
%%% (f) 記号論・LIS 側: Barthes(connotation の語の出所。ref.bib に Barthes1977
%%%     あり・一次未検証)/ Eco (1976) A Theory of Semiotics の unlimited semiosis
%%%     (記憶ベース・Web 未確認)/ §2.2 で引用済みの Mai 2001。
%%% "connotation" の語で立てるなら (f)、「経験の瞬間ごとに立ち上がる」の実質で
%%% 立てるなら (a)-(c) が近い。
What this paper offers is a part cut from that inexhaustible
whole---a \emph{principal connotation space}
(Section~\ref{subsec:bundle_reach})---and a place for each work
within it.

How objective, then, is this coordinate system? No criterion fixes
it from outside---had connotation an objective standard, it would
be denotation already. What fixes the cut is purpose: what a
library or an archive wants \emph{principal} to mean, what a music
service or a bookshop is run for. Each organization builds from
the content data it can gather, under a metric suited to its own
ends---an inductive bias at the meta level, standing before the
one of Section~\ref{subsec:two_genres}. And purpose shared is
compatibility gained: organizations that mean the same thing by
\emph{principal} can hold mutually compatible coordinates.

What this asks of the record is, in the end, one item: the content
data itself, kept in a form on which a locally Euclidean
metric---a standard inner product---is valid. Whoever holds that,
and shares the purpose, reproduces the space.
